\clearpage
\begin{landscape}
\begingroup
\fontsize{10}{13}\selectfont
\fancyhead[L]{\small NACHLASS. STÖRUNGEN DER PALLAS}
\renewcommand{\headrulewidth}{0pt}
\pdfbookmark[2]{Allgemeine Störungen, Erste Rechnung, [9]}{pallas-V-fragment9}
\begin{center}[9.]\end{center}
\par\noindent\hspace*{\parindent}[Es folgt nun die Berechnung der Störungen des halben Parameters \(p\); zu diesem Zweck ist nur die Integration \(\int Vn\,\mathrm{d}t\) auszuführen. Ganz in derselben Weise, wie bisher verfahren wurde, erhält man durch die Interpolationsrechnung \(V\) nach Vielfachen von \(M\) und \(M'\) entwickelt, und zwar in Bogensecunden ausgedrückt.
\par\noindent\hspace*{\parindent}Um die Störungen des Logarithmen des halben Parameters zu erhalten, ist der gefundene Ausdruck also noch durch \(206265\) zu dividiren; man erhält so für den \textsc{Brigg}ischen Logarithmen des halben Parameters in Einheiten der 7. Decimale:
\[
\frac{\mathrm{d}\log p}{n\,\mathrm{d}t}=-\frac{10000000}{206265}\,2\lambda\mathord{.}\,V,
\]
\noindent wo \(\lambda\) der Modul des \textsc{Brigg}ischen Logarithmensystems ist.
\par\noindent\hspace*{\parindent}Die Integration wird wie im Vorigen ausgeführt und es sei wieder das Beispiel für das Argument \(5M'-2M\) gegeben:
\begingroup
\newcommand{\VNinePoint}{\mathord{.}\mkern5mu}
\newcommand{\VNineFive}{\VNinePoint\VNinePoint\VNinePoint\VNinePoint\VNinePoint}
\begin{center}
\begin{tabular}{@{}l@{\quad}l@{}}
\(\log b\;\VNineFive\VNineFive\VNineFive\VNinePoint\VNinePoint0{,}75012\)&\\
\(\log b'\;\VNineFive\VNineFive\VNineFive\VNinePoint\VNinePoint9{,}58995_n\)&\\
\multicolumn{1}{r}{\(Q=93^0\,57'\,22''\)}&\\
\(\log\sin Q\;\VNineFive\VNineFive\VNinePoint\VNinePoint\VNinePoint\VNinePoint9{,}99897\)&\\
\(\log q\;\VNineFive\VNineFive\VNineFive\VNinePoint\VNinePoint0{,}75115\)&also das Störungsglied:\\
\(\log\left(5\frac{n'}{n}-2\right)\;\VNineFive\VNinePoint\VNinePoint\VNinePoint8{,}74966_n\)&\(-4226\cos(5M'-2M+93^0\,57'\,22'').\)\\
\cline{1-1}
\multicolumn{1}{r}{\(2{,}00149_n\)}&\\[3pt]
\(\log\frac{10000000}{206265}\,2\lambda\;\VNineFive\VNinePoint\VNinePoint1{,}62439\)&\\
\cline{1-1}
\multicolumn{1}{r}{\(3{,}62588_n\)}&
\end{tabular}
\end{center}
\endgroup
\par\noindent\hspace*{\parindent}Hiernach ergeben sich die Störungen des Logarithmen des halben Parameters, wie folgt, in Einheiten der 7. Decimale:
\begingroup
\renewcommand{\arraystretch}{1.15}
\begin{longtable}{@{}r@{}l@{}}
\(\delta\log p={}\)&\(+357\cos(M+68^0\,56'\,28'')\)\\
&\(+181\cos(2M+239\;50\;4)\)\\
&\(+59\cos(3M+296\;9\;41)\)\\[3pt]
&\(-39\cos(M'-4M+148^0\,54'\,52'')\)\\
&\(-100\cos(M'-3M+196\;53\;44)\)\\
&\(-226\cos(M'-2M+3\;29\;54)\)\\
&\(-269\cos(M'-M+162^0\,35'\,15'')\)\\
&\(+299\cos(M'+312\;40\;58)\)\\
&\(+110\cos(M'+M+115\;19\;51)\)\\
&\(+46\cos(M'+2M+226\;28\;46)\)\\
&\(+18\cos(M'+3M+359\;11\;31)\)\\[3pt]
&\(-26\cos(2M'-5M+35^0\,25'\,24'')\)\\
\SourceLeaf{528}&\(-26\cos(2M'-4M+65^0\,47'\,25'')\)\\
&\(-238\cos(2M'-3M+290\;11\;8)\)\\
&\(-1050\cos(2M'-2M+329\;46\;31)\)\\
&\(-4695\cos(2M'-M+144\;6\;20)\)\\
&\(+325\cos(2M'+319\;32\;42)\)\\
&\(+43\cos(2M'+M+107\;36\;3)\)\\
&\(+24\cos(2M'+2M+259\;14\;9)\)\\[3pt]
&\(-15\cos(3M'-6M+286^0\,28'\,43'))\)\\
&\(-5\cos(3M'-5M+244\;9\;21)\)\\
&\(-134\cos(3M'-4M+185\;33\;5)\)\\
&\(-324\cos(3M'-3M+235\;32\;0)\)\\
&\(-1137\cos(3M'-2M+23\;17\;27)\)\\
&\(+2848\cos(3M'-M+175\;0\;59)\)\\
&\(+120\cos(3M'+323\;12\;44)\)\\
&\(+23\cos(3M'+M+130\;24\;46)\)\\[3pt]
&\(-65\cos(4M'-5M+81^0\,18'\;6'')\)\\
&\(-107\cos(4M'-4M+154\;16\;43)\)\\
&\(-423\cos(4M'-3M+278\;56\;27)\)\\
&\(-1072\cos(4M'-2M+62\;26\;37)\)\\
&\(+394\cos(4M'-M+201\;46\;43)\)\\
&\(+46\cos(4M'+347\;40\;54)\)\\[3pt]
&\(-28\cos(5M'-6M+340^0\,21'\,57'')\)\\
&\(-48\cos(5M'-5M+80\;34\;4)\)\\
&\(-167\cos(5M'-4M+179\;55\;47)\)\\
&\(-356\cos(5M'-3M+317\;27\;26)\)\\
&\(-4226\cos(5M'-2M+93\;57\;22)\)\\
&\(+110\cos(5M'-M+230\;12\;34)\)\\
&\(+18\cos(5M'+10\;48\;14)\)\\[3pt]
&\(-26\cos(6M'-6M+358^0\,12'\;4'')\)\\
&\(-67\cos(6M'-5M+86\;40\;29)\)\\
&\(-152\cos(6M'-4M+215\;53\;22)\)\\
&\(-311\cos(6M'-3M+350^0\,9'\;1'')\)\\
&\(+353\cos(6M'-2M+122\;43\;27)\)\\
&\(+37\cos(6M'-M+257\;22\;35)\)\\[3pt]
&\(-14\cos(7M'-7M+267^0\,52'\;6'')\)\\
&\(-29\cos(7M'-6M+358\;30\;24)\)\\
&\(-68\cos(7M'-5M+117\;35\;50)\)\\
&\(-123\cos(7M'-4M+248\;40\;20)\)\\
&\(-395\cos(7M'-3M+20\;50\;27)\)\\
&\(+80\cos(7M'-2M+151\;44\;45)\)\\
&\(+14\cos(7M'-M+281\;40\;23)\)\\[3pt]
&\(-31\cos(8M'-6M+22^0\,30'\;5'')\)\\
&\(-57\cos(8M'-5M+149\;12\;21)\)\\
&\(-101\cos(8M'-4M+279\;34\;22)\)\\
&\(+517\cos(8M'-3M+50\;0\;55)\)\\
&\(+26\cos(8M'-2M+178\;4\;57)\)\\[3pt]
&\(-14\cos(9M'-7M+290^0\,38'\,30'')\)\\
&\(-27\cos(9M'-6M+51\;37\;9)\)\\
&\(-44\cos(9M'-5M+179\;33\;47)\)\\
&\(-99\cos(9M'-4M+308\;55\;32)\)\\
&\(+59\cos(9M'-3M+77\;10\;14)\)\\[3pt]
&\(-21\cos(10M'-6M+85^0\,3'\,29'')\)\\
&\(-35\cos(10M'-5M+209\;13\;29)\)\\
&\(-232\cos(10M'-4M+336\;0\;0)\)\\
&\(+17\cos(10M'-3M+103\;50\;34)\)\\[3pt]
&\(-16\cos(11M'-6M+110^0\,49'\,58'')\)\\
&\(-31\cos(11M'-5M+233\;0\;2)\)\\
&\(+52\cos(11M'-4M+3\;2\;22)\)
\end{longtable}
\begin{center}
\begin{tabular}{rl}
Säcular: jährlich&\(+117{,}48\)\\
täglich&\(+0{,}32163\).
\end{tabular}
\end{center}
\endgroup
\par\noindent\hspace*{\parindent}Aus dem Werthe des halben Parameters lässt sich auch die Excentricität finden, sobald die halbe grosse Axe bekannt ist; denn es ist \(p=a\cos\varphi^2\). Die Säcularänderung von \(\varphi\) kann schon hier angegeben werden; sie ist jährlich: \(-11''{,}13\), täglich \(-0''{,}03047\).]
\clearpage
\endgroup
\end{landscape}
